\begin{frame}[allowframebreaks]{Threshold Secret Sharing}
	\textbf{\textcolor{blue}{Shamir's \(t, n\) Threshold Scheme}}
	\begin{itemize}
		\item Is based on \textbf{threshold} access structure \(t\)-out-of-\(n\):
			\[
			\mathcal{A}_t = \{A\subseteq \mathcal{P}\;:\;|A|\ge k\}.
			\]
		\item Work over a finite field \(\mathbb{F}_q\) with \(q>n\)
		\item Choose a random polynomial of degree \(t-1\):
		\[
		P(x) = s + \alpha_1x + \cdots + \alpha_{t-1}x^{t-1},
		\quad \alpha_i \in \mathbb{F}_q.
		\]
		\item Give share \(P(i)\) to party \(P_i\).
		\item Any \(t\) shares interpolate \(P(x)\) and recover \(s=P(0)\); any \(t-1\) shares reveal nothing.
	\end{itemize}
	
	\textbf{\textcolor{blue}{Why Shamir is a Baseline}}
	\begin{itemize}
		\item Perfect privacy and correctness for threshold structures.
		\item Ideal information rate for threshold access structures \cite{brickell1989some}.
		\item Linearity enables efficient protocols and MPC constructions.
		\item Limitation: primarily threshold policies unless generalized.
	\end{itemize}
	
	\textbf{\textcolor{blue}{The need for a General Access Structure}}
	\begin{itemize}
		\item Real policies are not always “any \(t\) out of \(n\)”.
		\item General secret sharing schemes exist for any monotone access structure, but efficiency is the main challenge \cite{beimel2011secret}.
	\end{itemize}
\end{frame}

\begin{frame}{Benaloh-Leichter Approach}
	\begin{itemize}
		\item Represent an access policy as a monotone formula over parties.
		\item If an access structure can be described by a \textit{small monotone formula} \(\rightarrow\) efficient perfect secret sharing-scheme.
		\item For \(F = F_1 \vee \cdots \vee F_k\): replicate the secret down each branch.
		\item For \(F = F_1 \wedge \cdots \wedge F_k\): split the secret into \(k\) parts using a \((k,k)\) scheme, then recurse.
		\item Formula size directly affects share size blow-up.
	\end{itemize}
\end{frame}

\begin{frame}{Efficiency Challenge for General Access Structures}
	\begin{itemize}
		\item General constructions can have exponential share size in the worst case.
		\item There is a significant gap between known upper bounds and lower bounds for some families.
		\item Lower bounds and structural results for linear and multilinear schemes are surveyed in.
	\end{itemize}
\end{frame}

\begin{frame}{Ideal Secret Sharing and Structural Characterizations - \cite{brickell1989some}}
	\begin{itemize}
		\item Ideal scheme: share size equals secret size.
		\item Threshold structures are ideal; some non-threshold structures are also ideal.
		\item Links between ideality and combinatorial structures (e.g., matroids) appear in the classical literature.
	\end{itemize}
\end{frame}
